\documentclass[11pt,a4paper]{article}
\usepackage[T1]{fontenc}
\usepackage{lmodern}
\usepackage{amsmath,amssymb,amsthm,mathtools}
\usepackage[margin=27mm]{geometry}
\usepackage[hidelinks]{hyperref}

\DeclareMathOperator{\Tr}{Tr}
\DeclareMathOperator{\RePart}{Re}
\DeclareMathOperator{\ImPart}{Im}
\newcommand{\R}{\mathbb R}
\newcommand{\C}{\mathbb C}
\newcommand{\BZ}{\mathcal B}
\newcommand{\Torus}{\mathbb T^*}
\newcommand{\cH}{\mathcal H}
\newcommand{\ind}{\mathbf 1}
\newcommand{\e}{\mathrm e}
\newcommand{\trv}{\operatorname{tr}_{\rm vol}}
\newtheorem{theorem}{Theorem}
\newtheorem{lemma}{Lemma}
\newtheorem{remark}{Remark}
\setlength{\parindent}{0pt}
\setlength{\parskip}{0.45em}

\title{Residual conductivity at isolated conical crossings}
\author{}
\date{}

\begin{document}
\maketitle

\begin{abstract}
We prove a linear-response formula for a two-dimensional periodic Schr\"odinger
operator whose Fermi level meets the spectrum only at finitely many isolated
isotropic conical crossings. For a constant electric field switched on at time
zero, the long-time Ces\`aro mean of the first-order current exists. Each
crossing contributes one sixteenth of the identity conductivity tensor, in
atomic units and per spinless cone. The proof reduces the response to a
Bloch-fiber formula and evaluates its singular part using the two-band
effective Hamiltonian.
\end{abstract}

\section{Setting and statement}

Let $\mathcal R\subset\R^2$ be a lattice, let $\Omega$ be a fundamental cell,
and let
\[
 \mathcal R^*=\{K\in\R^2:K\cdot R\in2\pi\mathbb Z\text{ for all }R\in\mathcal R\}
\]
be its reciprocal lattice. We identify the Brillouin zone $\BZ$ with a
fundamental cell of $\R^2/\mathcal R^*$ and write $|\BZ|$ for its area. Thus
\begin{equation}\label{eq:area}
 |\Omega|\,|\BZ|=(2\pi)^2.
\end{equation}
Let $V\in H^1_{\rm per}(\Omega;\R)$ and consider
\[
 H=-\tfrac12\Delta+V\quad\text{on }L^2(\R^2),\qquad
 H_k=\tfrac12(-i\nabla+k)^2+V\quad\text{on }\cH:=L^2_{\rm per}(\Omega).
\]
The fiber operators have compact resolvent and common domain
$H^2_{\rm per}(\Omega)$. Their eigenvalues, counted with multiplicity, are
denoted by $\lambda_1(k)\leq\lambda_2(k)\leq\cdots$; set
$\lambda_0(k)=-\infty$.

Indeed, in two dimensions $V\in H^1(\Omega)\subset L^2(\Omega)$ and
$H^2(\Omega)\hookrightarrow L^\infty(\Omega)$, so multiplication by $V$ is
infinitesimally $-\Delta$-bounded on the fiber domain. The Kato--Rellich
theorem gives the stated self-adjoint realizations (Reed--Simon, Vol.~II,
Theorem X.12). The field Hamiltonian is taken with the self-adjoint Stark
closure specified in the hypotheses.

Fix $\mu_{\rm F}\in\R$. We assume that there are an integer $N\geq1$ and
distinct points $k_1,\ldots,k_M$ of $\R^2/\mathcal R^*$ such that
\begin{align}
 \lambda_N(k)&<\mu_{\rm F}<\lambda_{N+1}(k)
 &&(k\notin\{k_1,\ldots,k_M\}),\label{eq:gap}\\
 \lambda_{N-1}(k_i)&<\lambda_N(k_i)=\mu_{\rm F}
 =\lambda_{N+1}(k_i)<\lambda_{N+2}(k_i),\label{eq:cross}
\end{align}
and, for each $i$, there is $v_i>0$ for which
\begin{align}
 \lambda_N(k_i+q)&=\mu_{\rm F}-v_i|q|+O(|q|^2),\notag\\
 \lambda_{N+1}(k_i+q)&=\mu_{\rm F}+v_i|q|+O(|q|^2)
 \qquad(q\longrightarrow0).\label{eq:cone}
\end{align}
All points and neighborhoods below are understood on the reciprocal torus.

For real directions $e_\alpha,e_\beta\in\R^2$, let
\[
 J_\alpha=(-i\nabla)\cdot e_\alpha,\qquad X_\beta=x\cdot e_\beta,
 \qquad \gamma_0=\ind_{(-\infty,\mu_{\rm F}]}(H).
\]
For $\varepsilon\geq0$, set $H_\varepsilon=H-\varepsilon X_\beta$ (the
self-adjoint Stark realization) and
\[
 \gamma_\varepsilon(t)=\e^{-itH_\varepsilon}\gamma_0\e^{itH_\varepsilon}.
\]
The current density is the local trace per unit area
\[
 j_\varepsilon(t)=\frac1{|\Omega|}
 \Tr\bigl(\ind_\Omega J_\alpha\gamma_\varepsilon(t)\ind_\Omega\bigr),
\]
understood by the usual local trace interpretation when the field is present.
The Bloch evolution below shows that this quantity is well defined for finite
$t$ and that the evolved density remains lattice-periodic.

\begin{theorem}\label{thm:main}
Under assumptions \eqref{eq:gap}--\eqref{eq:cone}, the iterated limit exists
and is
\begin{equation}\label{eq:main}
 \lim_{T\to\infty}\frac1T\int_0^T
 \left(\lim_{\varepsilon\downarrow0}
       \frac{j_\varepsilon(t)}{\varepsilon}\right)dt
 =\frac{M}{16}\,e_\alpha\cdot e_\beta.
\end{equation}
\end{theorem}

The formula is for spinless particles. Each point $k_i$ is counted once, and
the conductivity tensor contributed by a single crossing is $\frac1{16}I_2$.

\section{Bloch-space linear response}

We use the Bloch transform with fibers $\cH$. In this representation,
$H$ is multiplication by $H_k$, $J_\alpha$ is multiplication by
\[
 J_\alpha(k)=\partial_{e_\alpha}H_k=(-i\nabla+k)\cdot e_\alpha,
\]
and $X_\beta$ acts as $i\partial_{e_\beta}$, with the standard reciprocal
boundary identifications. Since $V$ is real, complex conjugation identifies
the fibers at $k$ and $-k$.

Write $P(k)=\ind_{(-\infty,\mu_{\rm F}]}(H_k)$. Away from the crossing points,
$P(k)$ has rank $N$. Its value at the finite exceptional set is immaterial to
all integrals. For a periodic operator with fiber representative $A(k)$,
the trace-per-area normalization is
\begin{equation}\label{eq:blochtrace}
 \trv(A)=\frac{1}{|\Omega|\,|\BZ|}
       \int_\BZ\Tr_{\cH}(A(k))\,dk
       =\frac1{4\pi^2}\int_\BZ\Tr_{\cH}(A(k))\,dk.
\end{equation}

\begin{lemma}[Linear-response identity]\label{lem:response}
For each finite $t\geq0$, the right derivative
\[
 r_{\alpha\beta}(t):=
 \lim_{\varepsilon\downarrow0}\frac{j_\varepsilon(t)}{\varepsilon}
\]
exists. At points $k$ away from the crossings choose fiber eigenvectors
$u_{n,k}$ and put, for $n\leq N<m$,
\[
 \Delta_{mn}(k)=\lambda_m(k)-\lambda_n(k),\qquad
 A_{nm}^{\alpha\beta}(k)=
 \langle u_{n,k},J_\alpha(k)u_{m,k}\rangle
 \langle u_{m,k},J_\beta(k)u_{n,k}\rangle.
\]
Then, with the sum interpreted invariantly over occupied and unoccupied
spectral subspaces if there are degeneracies away from the Fermi level,
\begin{equation}\label{eq:response}
 \begin{aligned}
 r_{\alpha\beta}(t)=\frac1{4\pi^2}\int_\BZ
 \sum_{n\leq N<m}\frac{2}{\Delta_{mn}(k)^2}
 \Bigl[&\RePart A_{nm}^{\alpha\beta}(k)\sin(\Delta_{mn}(k)t)\\
 &+\ImPart A_{nm}^{\alpha\beta}(k)
       \bigl(1-\cos(\Delta_{mn}(k)t)\bigr)\Bigr]dk.
 \end{aligned}
\end{equation}
\end{lemma}

\begin{proof}
In Bloch variables the Stark Hamiltonian is $H_k-i\varepsilon
\partial_{e_\beta}$. Its Liouville equation for the fiber density is
\[
 \partial_t\Gamma_\varepsilon(t,k)
 +\varepsilon\partial_{e_\beta}\Gamma_\varepsilon(t,k)
 =-i[H_k,\Gamma_\varepsilon(t,k)],\qquad
 \Gamma_\varepsilon(0,k)=P(k).
\]
The characteristic flow in $k$ shows that $\Gamma_\varepsilon(t,\cdot)$ is a
rank-$N$ fiber density almost everywhere and hence remains periodic in real
space. At $\varepsilon=0$, $\Gamma_0(t,k)=P(k)$. Differentiating the transport
equation at zero gives, in the weak sense,
\[
 \partial_t D(t,k)+\partial_{e_\beta}P(k)=-i[H_k,D(t,k)],
 \qquad D(0,k)=0,
\]
so
\begin{equation}\label{eq:D}
 D(t,k)=-\int_0^t \e^{-i(t-s)H_k}
       (\partial_{e_\beta}P)(k)\e^{i(t-s)H_k}\,ds.
\end{equation}
The first derivative is legitimate after integration over $k$: near a cone
$P(k_i+q)$ has derivative $O(|q|^{-1})$, which is locally integrable in two
dimensions, and the finite-time oscillatory factor in \eqref{eq:D} gives the
same local integrability for the current trace. Equivalently, one may first
delete disks of radius $\delta$ around the crossings, differentiate on the
remaining compact set, and then let $\delta\downarrow0$ using these bounds.
Also, the equilibrium current vanishes: complex conjugation sends the fiber
at $k$ to that at $-k$, preserves $P$, and changes the sign of $J_\alpha$.
Thus $j_0(t)=0$, so the derivative just computed is the stated limit of
$j_\varepsilon(t)/\varepsilon$.

Differentiate $[H_k,P(k)]=0$. If $n$ is occupied, $m$ is unoccupied, and
$\Delta=\lambda_m-\lambda_n$, then
\[
 (\partial_{e_\beta}P)_{mn}
 =-\frac{(J_\beta)_{mn}}{\Delta},\qquad
 (\partial_{e_\beta}P)_{nm}
 =-\frac{(J_\beta)_{nm}}{\Delta}.
\]
There are no occupied-to-occupied or unoccupied-to-unoccupied terms in
$\partial P$. Inserting these expressions into \eqref{eq:D}, taking the fiber
trace with $J_\alpha(k)$, and integrating the two conjugate terms in $s$
gives
\[
 \frac{2}{\Delta^2}\left(
 \RePart A_{nm}^{\alpha\beta}\sin(\Delta t)
 +\ImPart A_{nm}^{\alpha\beta}(1-\cos(\Delta t))\right).
\]
Finally use \eqref{eq:blochtrace}. The sums converge locally uniformly away
from the crossing points; near each crossing only the touching pair is
singular, and its finite-time contribution is integrable by the same
$|q|^{-1}$ estimate. This proves the formula.
\end{proof}

\begin{lemma}[Time-reversal cancellation]\label{lem:TR}
After integration over $\BZ$, the imaginary parts of the summands in
\eqref{eq:response} vanish. Consequently,
\begin{equation}\label{eq:meanresponse}
 \frac1T\int_0^T r_{\alpha\beta}(t)\,dt
 =\frac1{4\pi^2}\int_\BZ\sum_{n\leq N<m}
 2\RePart A_{nm}^{\alpha\beta}(k)
 \frac{1-\cos(\Delta_{mn}(k)T)}{\Delta_{mn}(k)^3T}\,dk.
\end{equation}
\end{lemma}

\begin{proof}
Let $\mathcal C$ be complex conjugation on periodic cell functions. It obeys
$\mathcal C H_k\mathcal C=H_{-k}$ and
$\mathcal C J_\alpha(k)\mathcal C=-J_\alpha(-k)$. Pairing eigenvectors at
$k$ and $-k$ therefore gives
\[
 \Delta_{mn}(-k)=\Delta_{mn}(k),\qquad
 A_{nm}^{\alpha\beta}(-k)
 =\overline{A_{nm}^{\alpha\beta}(k)}.
\]
Thus the imaginary part is odd on the reciprocal torus and integrates to zero.
The average of $\sin(\Delta t)$ over $[0,T]$ is
$(1-\cos(\Delta T))/(\Delta T)$, which yields
\eqref{eq:meanresponse}.
\end{proof}

\section{The contribution of a conical crossing}

\begin{lemma}[Local two-band structure]\label{lem:local}
Fix one crossing $k_i$ with slope $v=v_i$. In a sufficiently small disk
$q\mapsto k_i+q$, the touching pair admits a smooth two-dimensional spectral
subspace and a smooth orthonormal frame in which its effective Hamiltonian is
\[
 H_{\rm eff}(q)=\mu_{\rm F}I+v\,Uq\cdot\boldsymbol\sigma+O(|q|^2),
 \qquad U^TU=I_2.
\]
Here $U:\R^2\to\R^3$ is a linear isometric embedding and
$\boldsymbol\sigma=(\sigma_1,\sigma_2,\sigma_3)$ are the Pauli matrices.
For the touching valence/conduction pair, uniformly in $\widehat q=q/|q|$,
\begin{align}
 \Delta(q)&=2v|q|+O(|q|^2),\label{eq:gaplocal}\\
 \RePart A^{\alpha\beta}(q)
 &=v^2\left(e_\alpha\cdot e_\beta
 -(e_\alpha\cdot\widehat q)(e_\beta\cdot\widehat q)\right)
 +O(|q|).\label{eq:Alocal}
\end{align}
\end{lemma}

\begin{proof}
The pair is isolated from the other fiber eigenvalues at $q=0$ by
\eqref{eq:cross}. The fibers form an analytic family in $k$: multiplication
by $V\in H^1(\Omega)\subset L^2(\Omega)$ is infinitesimally
$-\Delta$-bounded on $H^2(\Omega)$, and the $k$-dependence is polynomial.
The standard analytic spectral-projection and finite-dimensional reduction
theorems therefore give a smooth rank-two spectral projection and frame
(Reed--Simon, Vol.~IV, Theorems XII.5 and XII.12--13).

Write the linear Taylor part of the reduced Hermitian matrix as
$a(q)I+d(q)\cdot\boldsymbol\sigma$, with $a$ scalar and $d$ a linear map
from $\R^2$ to $\R^3$. Its eigenvalues are
$\mu_{\rm F}+a(q)\pm|d(q)|+O(|q|^2)$. The two expansions
\eqref{eq:cone} imply $a(q)=0$ and $|d(q)|=v|q|$. Hence
$d(q)=vUq$ for an isometric embedding $U$, proving the stated form.

Let $n(q)=U\widehat q+O(|q|)$ be the unit vector in the direction of the
traceless part. Differentiating the effective Hamiltonian in a direction
$e\in\R^2$ gives $v(Ue)\cdot\boldsymbol\sigma+O(|q|)$. The derivative of
the moving frame contributes only a commutator with
$H_{\rm eff}-\mu_{\rm F}I=O(|q|)$, so the physical current matrix has the
same leading term. For Pauli matrices and eigenvectors of
$n\cdot\boldsymbol\sigma$, the real part of the product of two interband
matrix elements of $a\cdot\boldsymbol\sigma$ and $b\cdot\boldsymbol\sigma$
is $a\cdot b-(a\cdot n)(b\cdot n)$. Taking
$a=vUe_\alpha$, $b=vUe_\beta$ proves \eqref{eq:Alocal}; the eigenvalue gap
gives \eqref{eq:gaplocal}.
\end{proof}

\section{Proof of the theorem}

\begin{proof}[Proof of Theorem~\ref{thm:main}]
By Lemma~\ref{lem:response}, the inner limit in \eqref{eq:main} is
$r_{\alpha\beta}(t)$. By Lemma~\ref{lem:TR}, its time average is
\eqref{eq:meanresponse}. We evaluate that expression as $T\to\infty$.

Choose disjoint disks around the $k_i$. On the complement, the direct gap
$\lambda_{N+1}-\lambda_N$ has a positive lower bound. All occupied/unoccupied
energy differences are then bounded below. The same is true in each disk for
every transition except the touching pair, by \eqref{eq:cross}. These regular
terms vanish as $T\to\infty$: indeed
\[
 \left|\frac{1-\cos(\Delta T)}{\Delta^3T}\right|
 \leq\frac{2}{\Delta^3T},
\]
and the current matrix-element sums are integrable uniformly in $k$. To see
the latter, on any region with $\Delta\geq g>0$, Cauchy--Schwarz gives, for
each occupied eigenvector,
\[
 \sum_{m>N}\frac{|(J_\alpha)_{nm}(J_\beta)_{mn}|}{\Delta_{mn}^3}
 \leq g^{-3}\|J_\alpha(k)u_{n,k}\|\,
                    \|J_\beta(k)u_{n,k}\|.
\]
There are only $N$ occupied bands and these norms are uniformly bounded on
the compact Brillouin torus. The same estimate applies to the nonsingular
transitions within the cone disks.

It remains to compute the touching pair near one $k_i$. Put
$q=r\widehat q$ and abbreviate
\[
 S(\widehat q)=e_\alpha\cdot e_\beta
 -(e_\alpha\cdot\widehat q)(e_\beta\cdot\widehat q).
\]
By Lemma~\ref{lem:local}, its unnormalized contribution to the integral in
\eqref{eq:meanresponse} is
\[
 I_i(T)=\int_{|q|<\delta}
 2\RePart A^{\alpha\beta}(q)
 \frac{1-\cos(\Delta(q)T)}{\Delta(q)^3T}\,d^2q.
\]
For fixed $\widehat q$, the function $r\mapsto s=\Delta(r\widehat q)$ is
strictly increasing for small $r$, with
$s=2vr+O(r^2)$ and $r\,dr=(s/(4v^2)+O(s^2))\,ds$. Combining this with
\eqref{eq:Alocal}, the radial integrand becomes
\[
 \left(\frac{S(\widehat q)}{2s^2T}+O\!\left(\frac1{sT}\right)\right)
 (1-\cos(sT))\,ds.
\]
The integrated error is $O((\log T)/T)$, uniformly in angle. The elementary
Dirichlet integral
\[
 \lim_{T\to\infty}\frac1T\int_0^{s_0}
       \frac{1-\cos(sT)}{s^2}\,ds=\frac\pi2
\]
therefore gives
\begin{align*}
 \lim_{T\to\infty} I_i(T)
 &=\frac\pi4\int_0^{2\pi}S(\widehat q(\theta))\,d\theta\\
 &=\frac\pi4\left(2\pi\,e_\alpha\cdot e_\beta
       -\pi\,e_\alpha\cdot e_\beta\right)
 =\frac{\pi^2}{4}\,e_\alpha\cdot e_\beta.
\end{align*}
Here we used
$\int_0^{2\pi}\widehat q_a(\theta)\widehat q_b(\theta)\,d\theta
=\pi\delta_{ab}$. Finally, the trace-per-area factor in
\eqref{eq:meanresponse} is $1/(4\pi^2)$ by \eqref{eq:area}. Thus one cone
contributes $(e_\alpha\cdot e_\beta)/16$. Summing the $M$ disjoint cone
contributions proves \eqref{eq:main}.
\end{proof}

\begin{thebibliography}{9}
\bibitem{RS2}
M. Reed and B. Simon, \emph{Methods of Modern Mathematical Physics, Vol.~II:
Fourier Analysis, Self-Adjointness}, Academic Press, New York, 1975.

\bibitem{RS4}
M. Reed and B. Simon, \emph{Methods of Modern Mathematical Physics, Vol.~IV:
Analysis of Operators}, Academic Press, New York, 1978.
\end{thebibliography}

\end{document}
